%% linguex's \ExLBr / \ExRBr, which a document written for linguex uses to
%% get [1] instead of (1).  They were accepted and ignored here until the
%% characters moved into them: the document compiled and only the page was
%% wrong, which is the one kind of incompatibility a drop-in replacement
%% cannot have.
%%
%% Three things at once: the main pair is honoured; the FOOTNOTE pair is
%% separate, so a footnote example still prints (i) when only \ExLBr moved
%% (as in linguex, where they are four commands and not two); and
%% \theExLBr still overrides both, since it is the layer that also carries
%% the parenthesis-suppression switch.
\input{_preamble}
\renewcommand{\ExLBr}{[}
\renewcommand{\ExRBr}{]}
\begin{document}
\ex.\label{ex:one} EXLBRMAIN the first example.

\ex. EXLBRSECOND the second.

EXLBRREF \ref{ex:one} is how it is referred to.

EXLBRFNHOST a host sentence.\footnote{EXLBRFNTEXT \ex. EXLBRFN in a footnote.}

%% Doubled brackets rather than angles: pdftotext -bbox reports its
%% words as XML and would hand the suite "&lt;3&gt;" to assert on.
{\renewcommand{\theExLBr}{[[}\renewcommand{\theExRBr}{]]}%
\ex. EXLBROVERRIDE the layer above still wins.

}
\end{document}
